\section{Introduction}

Device fingerprinting is an essential step in cyber reconnaissance, letting adversaries reveal unique
features of a target network and design stealthy attack strategies. These techniques are increasingly
critical in industrial control systems (ICSs) such as power grids, where adversaries use IP-based
control networks as entry points to inflict physical damage. Consequently, fingerprinting shifts the
focus from identifying visited websites and user behavior to revealing device models and control
operations that are critical to ICS attacks. In the 2015 attack that disrupted the Ukrainian power
grid and the Stuxnet attack that disrupted Iranian nuclear facilities, it is widely believed that
adversaries stayed in the target systems for at least six months to perform cyber
reconnaissance~\cite{leeAnalysisCyberAttack2016}.

Power grids carry their measurements and control commands over protocols such as DNP3 (Distributed
Network Protocol 3)~\cite{ieeeIEEEStandardElectric2012}, which was deployed with no encryption by
default~\cite{eastTaxonomyAttacksDNP32009}. Consequently, any observer on the communication path can read the
exchange, and need not even decode the payload. Because a control-related attack, such as a false-data
injection attack on state estimation, works only when the attacker knows the target
device~\cite{liuFalseDataInjection2009}, the attacker first learns which devices communicate and how they behave.
For this step, the response timing leaks enough on its own~\cite{formbyWhosControlYour2016}: from it alone, an
observer can identify a device, its type, and its model~\cite{kohnoRemotePhysicalDevice2005, guFingerprintingCyberPhysicalSystem2018}.

To disrupt device fingerprinting, many studies present network traffic
obfuscation~\cite{wrightTrafficMorphingEfficient2009, dyerPeekaBooStillSee2012}. These defenses target the network-level features that
fingerprinting relies on~\cite{sirinamDeepFingerprintingUndermining2018}, in two directions. First, padding and splitting change the
distribution of packet sizes~\cite{wrightTrafficMorphingEfficient2009}. Second, delaying packets and adding dummy ones change the
inter-packet timing~\cite{dyerPeekaBooStillSee2012, caiSystematicApproachDeveloping2014}. Since manipulating traffic at the host adds runtime overhead,
recent work offloads obfuscation onto programmable switches, which reshape communication at much
higher line rates than CPUs~\cite{meierDittoWANTraffic2022, wangMinosLightweightDynamic2025, xieSecuritasDefendingTraffic2026}.

Unfortunately, these methods do not transfer to device fingerprinting in ICS environments, for two
reasons. First, the leaked features differ. General obfuscation targets flow-level size and timing
patterns~\cite{meierDittoWANTraffic2022}, whereas an ICS device fingerprint is carried by the response timing of the
protocol exchange, which stays visible even when the payload is encrypted~\cite{formbyWhosControlYour2016}. Second, ICS
protocols leave little room to add bytes or fabricate traffic. Because DNP3 validates its payload,
padding and injected packets break the exchange: the master reassembles each response and checks a
cyclic redundancy check for every sixteen application bytes~\cite{ieeeIEEEStandardElectric2012}. Consequently, a defense for
this setting must hide timing that stays visible under encryption while touching no byte.

In this paper, we present \sysname, an in-network defense that meets both conditions. \sysname
normalizes an outstation's fingerprintable timing in the data plane of a single programmable switch at
the outstation edge, without changing the endpoints and without modifying a DNP3 byte. For a read, it
holds the acknowledgment and the response in switch queues and releases them at fixed offsets from the
request, so the master-visible timing becomes a single policy value. For a control command, it
normalizes the timing on a second, independent scheduling lane, so the operation time also becomes a
single value. Because \sysname only holds and forwards packets, it changes no byte. We implement it as
one P4 program on a single switch pipeline and evaluate it on a physical SEL-751 relay.

We make the following contributions.

\begin{itemize}
  \item \textbf{Normalizes ICS device-fingerprinting timing in the data plane.} We propose an
    in-network framework that removes the timing features an on-path observer uses to fingerprint a
    DNP3 outstation, using one switch at the outstation edge and changing no endpoint and no DNP3
    byte. To the best of our knowledge, this is the first in-network defense that normalizes ICS
    device-fingerprinting features on real hardware.
  \item \textbf{Normalizes the cross-layer response time.} As a first case study, reproducing the
    read-side feature that Formby et al. use~\cite{formbyWhosControlYour2016}, we hold the acknowledgment and the
    response and release them at fixed offsets from the request, driving the master-visible timing to
    a single policy value.
  \item \textbf{Normalizes the control-command timing on an independent lane.} As a second case study,
    covering the physical-operation feature of the same attack~\cite{formbyWhosControlYour2016}, we hold the control
    response, pin the master-visible acknowledgment and echo to the request, and show why the read and
    control paths need two independent scheduling lanes.
  \item \textbf{Evaluates on a physical relay.} We evaluate \sysname on a physical SEL-751 with a
    reproducible capture-to-claim analysis, an explicit safety result, and a stated evidence boundary.
\end{itemize}
